欢迎来到视觉组。

作为比你们早走几步的学长，我们最想说的不是“这里有多难”，而是：别急。刚开始看不懂代码、配不好环境、被一行报错困住一下午，都是再正常不过的事。基础打得慢一点没关系，真正危险的是似懂非懂地往下冲，最后连问题出在哪里都说不清。

刚入组时，先用手边的电脑把一个小程序跑起来。第 1 章给出了 Windows 和 Ubuntu 两条路线，任选熟悉的系统和一套 IDE，就能观察一行代码执行前后的变化。等到卷一再补 Linux、C++ 和 CMake，你会知道那些命令和规则具体用在哪里。

AI 也可以从第一次练习就用起来。它能解释报错、提供思路，也可能误解题目，或给出一份能编译、结果却不对的修改。第 2 章从刚跑过的计分器入手，练习把问题说清楚、选对工具，再亲手核对结果。桌面或终端的操作遇到再查；要用工作树、撤回仓库改动或交付可复现的修改时，再用卷一的 Git 和 CMake 基础做项目实作。

程序到了另一台机器以后，还要核对版本、工作目录、依赖和日志。卷一第 4 章会接着带你走完这条实机路线。

遇到问题，先读报错、查资料、动手试，再带着现象和尝试过的方法来问。不要害怕求助，也不要把“求助”等同于把问题丢给别人。能把环境、操作、预期和结果讲清楚，是对自己负责，也是对队友时间的尊重。

写下的每一行代码，迟早都会由自己或队友来维护。动手之前先备份，修改之前先确认分支，提交之前再检查一遍；不确定会造成什么后果的命令，不要靠胆量执行。赛场不会因为一句“在我电脑上能跑”就网开一面，可靠往往来自那些看起来不起眼的习惯。

也别把一时的进度当成能力，更别拿自己的起点和别人的积累比较。承认不会，耐心补课，遇到问题不甩锅，答应的事情有回音——这些品质比偶尔写出一段漂亮代码更难得。

希望你们保留好奇心，也保留敬畏心。路还很长，慢慢来，但别停下。
